% Clase 14 — esfera de conductividad g_1 inmersa en un conductor g_2, con
% densidad de corriente uniforme lejos. Caso límite g_1 -> 0: la corriente
% le da la vuelta. Las líneas NO están dibujadas a ojo: salen de la solución
% exacta de Laplace con J_n = 0 en r = a,
%   psi = (1/2) J_0 rho^2 (1 - a^3/r^3),
% generadas con un script Python puro (bisección en r para cada theta).
% Parámetros: a = 1; líneas que lejos pasan a b = 0.3, 0.6, 0.95, 1.35, 1.75 del eje.
\begin{tikzpicture}[scale=1.05]

  \fill[figgray!10] (-3.2,-2.05) rectangle (3.2,2.05);
  \draw[figline] plot[smooth] coordinates {(-3.200,0.306) (-2.984,0.306) (-2.782,0.307) (-2.608,0.309) (-2.457,0.310) (-2.210,0.314) (-2.016,0.319) (-1.862,0.325) (-1.684,0.335) (-1.512,0.350) (-1.364,0.373) (-1.227,0.409) (-1.091,0.472) (-0.974,0.556) (-0.866,0.650) (-0.756,0.744) (-0.642,0.828) (-0.515,0.903) (-0.381,0.962) (-0.241,1.003) (-0.097,1.026) (0.049,1.029) (0.193,1.013) (0.335,0.978) (0.471,0.925) (0.601,0.855) (0.717,0.775) (0.827,0.684) (0.932,0.592) (1.045,0.502) (1.169,0.431) (1.318,0.383) (1.478,0.355) (1.635,0.339) (1.797,0.328) (2.016,0.319) (2.210,0.314) (2.457,0.310) (2.608,0.309) (2.782,0.307) (2.984,0.306) (3.200,0.306)};
  \draw[figvec] (-2.782,0.307) -- (-2.608,0.309);
  \draw[figline] plot[smooth] coordinates {(-3.200,-0.306) (-2.984,-0.306) (-2.782,-0.307) (-2.608,-0.309) (-2.457,-0.310) (-2.210,-0.314) (-2.016,-0.319) (-1.862,-0.325) (-1.684,-0.335) (-1.512,-0.350) (-1.364,-0.373) (-1.227,-0.409) (-1.091,-0.472) (-0.974,-0.556) (-0.866,-0.650) (-0.756,-0.744) (-0.642,-0.828) (-0.515,-0.903) (-0.381,-0.962) (-0.241,-1.003) (-0.097,-1.026) (0.049,-1.029) (0.193,-1.013) (0.335,-0.978) (0.471,-0.925) (0.601,-0.855) (0.717,-0.775) (0.827,-0.684) (0.932,-0.592) (1.045,-0.502) (1.169,-0.431) (1.318,-0.383) (1.478,-0.355) (1.635,-0.339) (1.797,-0.328) (2.016,-0.319) (2.210,-0.314) (2.457,-0.310) (2.608,-0.309) (2.782,-0.307) (2.984,-0.306) (3.200,-0.306)};
  \draw[figvec] (-2.782,-0.307) -- (-2.608,-0.309);
  \draw[figline] plot[smooth] coordinates {(-3.200,0.609) (-2.952,0.611) (-2.747,0.614) (-2.570,0.617) (-2.416,0.620) (-2.220,0.626) (-2.056,0.633) (-1.877,0.643) (-1.730,0.654) (-1.582,0.670) (-1.439,0.691) (-1.290,0.722) (-1.144,0.764) (-1.004,0.817) (-0.874,0.874) (-0.738,0.937) (-0.608,0.992) (-0.471,1.042) (-0.332,1.081) (-0.193,1.106) (-0.053,1.119) (0.088,1.117) (0.228,1.101) (0.367,1.072) (0.505,1.031) (0.643,0.978) (0.774,0.920) (0.903,0.861) (1.036,0.804) (1.170,0.756) (1.306,0.718) (1.460,0.687) (1.609,0.666) (1.764,0.651) (1.918,0.640) (2.108,0.630) (2.281,0.624) (2.490,0.619) (2.655,0.615) (2.845,0.613) (3.067,0.610) (3.200,0.609)};
  \draw[figvec] (-2.747,0.614) -- (-2.570,0.617);
  \draw[figline] plot[smooth] coordinates {(-3.200,-0.609) (-2.952,-0.611) (-2.747,-0.614) (-2.570,-0.617) (-2.416,-0.620) (-2.220,-0.626) (-2.056,-0.633) (-1.877,-0.643) (-1.730,-0.654) (-1.582,-0.670) (-1.439,-0.691) (-1.290,-0.722) (-1.144,-0.764) (-1.004,-0.817) (-0.874,-0.874) (-0.738,-0.937) (-0.608,-0.992) (-0.471,-1.042) (-0.332,-1.081) (-0.193,-1.106) (-0.053,-1.119) (0.088,-1.117) (0.228,-1.101) (0.367,-1.072) (0.505,-1.031) (0.643,-0.978) (0.774,-0.920) (0.903,-0.861) (1.036,-0.804) (1.170,-0.756) (1.306,-0.718) (1.460,-0.687) (1.609,-0.666) (1.764,-0.651) (1.918,-0.640) (2.108,-0.630) (2.281,-0.624) (2.490,-0.619) (2.655,-0.615) (2.845,-0.613) (3.067,-0.610) (3.200,-0.609)};
  \draw[figvec] (-2.747,-0.614) -- (-2.570,-0.617);
  \draw[figline] plot[smooth] coordinates {(-3.200,0.964) (-3.051,0.965) (-2.899,0.967) (-2.697,0.971) (-2.520,0.975) (-2.365,0.980) (-2.185,0.986) (-2.029,0.994) (-1.862,1.005) (-1.718,1.016) (-1.570,1.031) (-1.421,1.050) (-1.276,1.073) (-1.135,1.100) (-0.996,1.130) (-0.858,1.162) (-0.720,1.195) (-0.577,1.227) (-0.430,1.256) (-0.285,1.277) (-0.142,1.290) (0.000,1.294) (0.142,1.290) (0.285,1.277) (0.430,1.256) (0.577,1.227) (0.720,1.195) (0.858,1.162) (0.996,1.130) (1.135,1.100) (1.276,1.073) (1.421,1.050) (1.570,1.031) (1.718,1.016) (1.862,1.005) (2.029,0.994) (2.185,0.986) (2.365,0.980) (2.520,0.975) (2.697,0.971) (2.899,0.967) (3.051,0.965) (3.200,0.964)};
  \draw[figvec] (-2.899,0.967) -- (-2.697,0.971);
  \draw[figline] plot[smooth] coordinates {(-3.200,-0.964) (-3.051,-0.965) (-2.899,-0.967) (-2.697,-0.971) (-2.520,-0.975) (-2.365,-0.980) (-2.185,-0.986) (-2.029,-0.994) (-1.862,-1.005) (-1.718,-1.016) (-1.570,-1.031) (-1.421,-1.050) (-1.276,-1.073) (-1.135,-1.100) (-0.996,-1.130) (-0.858,-1.162) (-0.720,-1.195) (-0.577,-1.227) (-0.430,-1.256) (-0.285,-1.277) (-0.142,-1.290) (0.000,-1.294) (0.142,-1.290) (0.285,-1.277) (0.430,-1.256) (0.577,-1.227) (0.720,-1.195) (0.858,-1.162) (0.996,-1.130) (1.135,-1.100) (1.276,-1.073) (1.421,-1.050) (1.570,-1.031) (1.718,-1.016) (1.862,-1.005) (2.029,-0.994) (2.185,-0.986) (2.365,-0.980) (2.520,-0.975) (2.697,-0.971) (2.899,-0.967) (3.051,-0.965) (3.200,-0.964)};
  \draw[figvec] (-2.899,-0.967) -- (-2.697,-0.971);
  \draw[figline] plot[smooth] coordinates {(-3.200,1.367) (-3.031,1.369) (-2.857,1.372) (-2.699,1.375) (-2.555,1.378) (-2.382,1.384) (-2.227,1.389) (-2.055,1.396) (-1.901,1.404) (-1.737,1.414) (-1.591,1.425) (-1.438,1.438) (-1.281,1.454) (-1.140,1.469) (-0.993,1.486) (-0.843,1.505) (-0.701,1.521) (-0.554,1.538) (-0.411,1.551) (-0.260,1.561) (-0.111,1.567) (0.037,1.568) (0.185,1.565) (0.335,1.557) (0.475,1.545) (0.620,1.531) (0.771,1.513) (0.917,1.496) (1.057,1.479) (1.209,1.461) (1.358,1.446) (1.502,1.433) (1.662,1.420) (1.817,1.409) (1.961,1.401) (2.121,1.393) (2.264,1.388) (2.423,1.382) (2.601,1.377) (2.750,1.374) (2.913,1.371) (3.094,1.368) (3.200,1.367)};
  \draw[figvec] (-2.857,1.372) -- (-2.699,1.375);
  \draw[figline] plot[smooth] coordinates {(-3.200,-1.367) (-3.031,-1.369) (-2.857,-1.372) (-2.699,-1.375) (-2.555,-1.378) (-2.382,-1.384) (-2.227,-1.389) (-2.055,-1.396) (-1.901,-1.404) (-1.737,-1.414) (-1.591,-1.425) (-1.438,-1.438) (-1.281,-1.454) (-1.140,-1.469) (-0.993,-1.486) (-0.843,-1.505) (-0.701,-1.521) (-0.554,-1.538) (-0.411,-1.551) (-0.260,-1.561) (-0.111,-1.567) (0.037,-1.568) (0.185,-1.565) (0.335,-1.557) (0.475,-1.545) (0.620,-1.531) (0.771,-1.513) (0.917,-1.496) (1.057,-1.479) (1.209,-1.461) (1.358,-1.446) (1.502,-1.433) (1.662,-1.420) (1.817,-1.409) (1.961,-1.401) (2.121,-1.393) (2.264,-1.388) (2.423,-1.382) (2.601,-1.377) (2.750,-1.374) (2.913,-1.371) (3.094,-1.368) (3.200,-1.367)};
  \draw[figvec] (-2.857,-1.372) -- (-2.699,-1.375);
  \draw[figline] plot[smooth] coordinates {(-3.200,1.769) (-3.048,1.770) (-2.893,1.773) (-2.750,1.775) (-2.574,1.779) (-2.414,1.783) (-2.267,1.788) (-2.099,1.793) (-1.946,1.799) (-1.805,1.805) (-1.649,1.813) (-1.506,1.820) (-1.351,1.829) (-1.208,1.838) (-1.054,1.849) (-0.910,1.858) (-0.757,1.868) (-0.610,1.876) (-0.468,1.884) (-0.315,1.890) (-0.164,1.893) (-0.015,1.895) (0.134,1.894) (0.284,1.891) (0.437,1.885) (0.578,1.878) (0.723,1.870) (0.875,1.860) (1.018,1.851) (1.168,1.841) (1.309,1.832) (1.460,1.823) (1.600,1.815) (1.752,1.808) (1.917,1.800) (2.068,1.794) (2.233,1.789) (2.376,1.784) (2.533,1.780) (2.704,1.776) (2.893,1.773) (3.048,1.770) (3.200,1.769)};
  \draw[figvec] (-2.893,1.773) -- (-2.750,1.775);
  \draw[figline] plot[smooth] coordinates {(-3.200,-1.769) (-3.048,-1.770) (-2.893,-1.773) (-2.750,-1.775) (-2.574,-1.779) (-2.414,-1.783) (-2.267,-1.788) (-2.099,-1.793) (-1.946,-1.799) (-1.805,-1.805) (-1.649,-1.813) (-1.506,-1.820) (-1.351,-1.829) (-1.208,-1.838) (-1.054,-1.849) (-0.910,-1.858) (-0.757,-1.868) (-0.610,-1.876) (-0.468,-1.884) (-0.315,-1.890) (-0.164,-1.893) (-0.015,-1.895) (0.134,-1.894) (0.284,-1.891) (0.437,-1.885) (0.578,-1.878) (0.723,-1.870) (0.875,-1.860) (1.018,-1.851) (1.168,-1.841) (1.309,-1.832) (1.460,-1.823) (1.600,-1.815) (1.752,-1.808) (1.917,-1.800) (2.068,-1.794) (2.233,-1.789) (2.376,-1.784) (2.533,-1.780) (2.704,-1.776) (2.893,-1.773) (3.048,-1.770) (3.200,-1.769)};
  \draw[figvec] (-2.893,-1.773) -- (-2.750,-1.775);
  % línea del eje: llega al punto de estancamiento y sale por el otro lado
  \draw[figline] (-3.2,0) -- (-1,0);
  \draw[figline] (1,0) -- (3.2,0);
  \draw[figvec] (-3.0,0) -- (-2.6,0);

  % esfera
  \fill[white] (0,0) circle (1);
  \fill[figamber!18] (0,0) circle (1);
  \draw[figamber, line width=1pt] (0,0) circle (1);
  \node[figlbl] at (0,0) {$g_1 \to 0$};
  \node[figlbl, fill=figgray!10, inner sep=1.5pt] at (2.55,1.72) {$g_2$};
  \node[figlbl, fill=figgray!10, inner sep=1.5pt] at (-2.55,1.72) {$\vec J$};

\end{tikzpicture}
